\begin{theorem}\label{thm:fold-window6-chi-layered-spectrum-reconstruction-by-12-moments}
Assume \(m=6\). By Assumption \ref{cor:fold-bin-m6-fiber-hist},
\[
\begin{aligned}
d_6(x)&\in\{2,3,4\},\\
\#\{x\in X_6:d_6(x)=2\}&=8,\\
\#\{x\in X_6:d_6(x)=3\}&=4,\\
\#\{x\in X_6:d_6(x)=4\}&=9.
\end{aligned}
\]
and \(|X_6|=21\).
Write
\[
\chi(x)=(\chi_1(x),\chi_2(x))\in\{\pm1\}^2.
\]
For \(s,t\in\{0,1\}\) and \(q\ge 0\), define
\[
M_{s,t}(q)
:=
2^6\,\tau_{6,1}\!\left(
\varepsilon_{\mathrm{scan}}^{\,s}\varepsilon_{\mathrm{rev}}^{\,t}\,
\mathrm{IndW}(E_{6,1})^{\,q}
\right)
=
\sum_{x\in X_6}\chi_1(x)^s\chi_2(x)^t\,d_6(x)^{q+1}.
\]
For each \((\alpha,\beta)\in\{\pm1\}^2\) and \(j\in\{2,3,4\}\), let
\[
n_j^{\alpha,\beta}
:=
\#\{x\in X_6:\chi(x)=(\alpha,\beta),\ d_6(x)=j\}.
\]
Then the twelve numbers
\[
\{M_{s,t}(q):(s,t)\in\{0,1\}^2,\ q=0,1,2\}
\]
determine all twelve integers \(n_j^{\alpha,\beta}\). More precisely, if
\[
p_k^{\alpha,\beta}
:=
\sum_{\chi(x)=(\alpha,\beta)}d_6(x)^k
=
\frac14\sum_{s,t\in\{0,1\}}\alpha^s\beta^t\,M_{s,t}(k-1),
\qquad k=1,2,3,
\]
then
\[
n_2^{\alpha,\beta}
=
3p_1^{\alpha,\beta}-\frac74p_2^{\alpha,\beta}+\frac14p_3^{\alpha,\beta},
\]
\[
n_3^{\alpha,\beta}
=
-\frac83p_1^{\alpha,\beta}+2p_2^{\alpha,\beta}-\frac13p_3^{\alpha,\beta},
\]
\[
n_4^{\alpha,\beta}
=
\frac34p_1^{\alpha,\beta}-\frac58p_2^{\alpha,\beta}+\frac18p_3^{\alpha,\beta}.
\]
Consequently,
\[
Q_6^{\alpha,\beta}(z)
:=
\prod_{\chi(x)=(\alpha,\beta)}(1-d_6(x)z)
=
(1-2z)^{n_2^{\alpha,\beta}}(1-3z)^{n_3^{\alpha,\beta}}(1-4z)^{n_4^{\alpha,\beta}}.
\]
\end{theorem}

\begin{proof}
By Theorem \ref{thm:fold-groupoid-z2x2-joint-spectral-measure}, the quantities \(M_{s,t}(q)\) are the signed power sums of the \(\chi\)-layered fiber spectrum. Proposition \ref{prop:fold-chi-layered-fiber-spectrum-certified-by-signed-powersums} therefore gives
\[
p_k^{\alpha,\beta}
=
\frac14\sum_{s,t\in\{0,1\}}\alpha^s\beta^t\,M_{s,t}(k-1),
\qquad k=1,2,3.
\]

Since the only fiber sizes in window \(6\) are \(2\), \(3\), and \(4\), the multiplicities \(n_2^{\alpha,\beta}\), \(n_3^{\alpha,\beta}\), and \(n_4^{\alpha,\beta}\) satisfy
\[
p_1^{\alpha,\beta}=2n_2^{\alpha,\beta}+3n_3^{\alpha,\beta}+4n_4^{\alpha,\beta},
\]
\[
p_2^{\alpha,\beta}=4n_2^{\alpha,\beta}+9n_3^{\alpha,\beta}+16n_4^{\alpha,\beta},
\]
\[
p_3^{\alpha,\beta}=8n_2^{\alpha,\beta}+27n_3^{\alpha,\beta}+64n_4^{\alpha,\beta}.
\]
The coefficient matrix
\[
\begin{pmatrix}
2&3&4\\
4&9&16\\
8&27&64
\end{pmatrix}
\]
has determinant \(48\neq 0\), so the system is invertible. Subtracting twice the first equation from the second and twice the second equation from the third gives
\[
3n_3^{\alpha,\beta}+8n_4^{\alpha,\beta}=p_2^{\alpha,\beta}-2p_1^{\alpha,\beta},
\]
\[
9n_3^{\alpha,\beta}+32n_4^{\alpha,\beta}=p_3^{\alpha,\beta}-2p_2^{\alpha,\beta}.
\]
One further elimination step yields
\[
8n_4^{\alpha,\beta}=6p_1^{\alpha,\beta}-5p_2^{\alpha,\beta}+p_3^{\alpha,\beta},
\]
which is the stated formula for \(n_4^{\alpha,\beta}\). Substituting back into the previous relation gives the displayed formula for \(n_3^{\alpha,\beta}\), and then the first equation gives the formula for \(n_2^{\alpha,\beta}\).
\end{proof}

\begin{theorem}\label{thm:fold-window6-center-three-observables-dimension-defect}
Assume \(m=6\), and write
\[
A_{6,1}:=\CC[\mathcal R_6],
\qquad
Z_6:=Z(A_{6,1}),
\qquad
Z_6^{(3)}:=\CC\!\left[\mathrm{IndW}(E_{6,1}),\varepsilon_{\mathrm{scan}},\varepsilon_{\mathrm{rev}}\right]\subseteq Z_6.
\]
Then
\[
\dim_{\CC}Z_6^{(3)}\le 12.
\]
Since \(|X_6|=21\), it follows that
\[
\dim_{\CC}Z_6=21,
\qquad
\dim_{\CC}(Z_6/Z_6^{(3)})\ge 9.
\]
Equivalently, the three central observables \(\mathrm{IndW}(E_{6,1})\), \(\varepsilon_{\mathrm{scan}}\), and \(\varepsilon_{\mathrm{rev}}\) do not suffice to distinguish all \(21\) stable types in window \(6\).

Moreover,
\[
\dim_{\CC}Z_6^{(3)}=12
\iff
V_{\alpha,\beta}=\{2,3,4\}
\quad\text{for every }(\alpha,\beta)\in\{\pm1\}^2.
\]
\end{theorem}

\begin{proof}
By Assumption \ref{cor:fold-bin-m6-fiber-hist}, the possible fiber multiplicities in window \(6\) are exactly \(2,3,4\). Hence
\[
V_{\alpha,\beta}\subseteq\{2,3,4\},
\qquad
|V_{\alpha,\beta}|\le 3
\]
for every sector \((\alpha,\beta)\). Proposition \ref{prop:fold-center-resolution-criterion-indw-eps} therefore gives
\[
\dim_{\CC}Z_6^{(3)}
=
\sum_{(\alpha,\beta)\in\{\pm1\}^2}|V_{\alpha,\beta}|
\le 12.
\]

On the other hand, the Wedderburn decomposition
\[
A_{6,1}\cong \bigoplus_{x\in X_6}M_{d_6(x)}(\CC)
\]
shows that \(A_{6,1}\) has \(|X_6|=21\) simple blocks, hence \(\dim_{\CC}Z_6=21\). The lower bound on \(\dim_{\CC}(Z_6/Z_6^{(3)})\) follows at once.

Under the identification \(Z_6\cong \CC^{X_6}\), a unital subalgebra separates all points of \(X_6\) if and only if it is the whole function algebra: once the points are separated, Lagrange interpolation produces the characteristic functions of singletons. Since \(\dim_{\CC}Z_6^{(3)}<21=\dim_{\CC}Z_6\), the algebra generated by \(\mathrm{IndW}(E_{6,1})\), \(\varepsilon_{\mathrm{scan}}\), and \(\varepsilon_{\mathrm{rev}}\) cannot separate all stable types.

Finally, equality \(\dim_{\CC}Z_6^{(3)}=12\) holds if and only if \(|V_{\alpha,\beta}|=3\) in every sector, which is the same as requiring \(V_{\alpha,\beta}=\{2,3,4\}\) for all \((\alpha,\beta)\).
\end{proof}

\endinput
